%% =====================================================================
%%  T-04-03  Vanity as a Handle
%%  -------------------------------------------------------------------
%%  One idea per file. Core is mandatory. Slots in canonical order.
%%  Max 700 words. No layout commands. No filler. New source material
%%  for this entry goes in sources/<BOOK>/T-04-03.tex -- never by
%%  rewriting this file (docs/RULES.md R0.1).
%% =====================================================================
%% @kind: topic
%% @id: T-04-03
%% @title: Vanity as a Handle
%% @subtitle: Status sensitivity is the most widely distributed weakness
%% @chapter: c04
%% @order: 30
%% @status: stable
%% @sources: B3, B1
%% @updated: 2026-09-19

\topic{T-04-03}{Vanity as a Handle}{Status sensitivity is the most widely distributed weakness}

\begin{core}
Almost everyone has a rank in mind for themselves, and almost everyone suspects it is not the rank others assign them. That gap is the most commonly available handle there is, and it can be worked by flattering it, by threatening it, or by appearing to occupy it.
\end{core}
\begin{mechanism}
Status is the one good that cannot be self-supplied, because it consists entirely of other people's assessments. That makes the person permanently dependent on external confirmation, and therefore permanently responsive to whoever appears to be able to give or withhold it. Flattery works on the gap between self-rank and perceived rank. Provocation works by threatening the same gap. Superiority works by occupying a rank the other person wanted. The reason the handle is so reliable is that defending status feels like defending self, and people will trade material interests to do it.
\end{mechanism}
\begin{tells}
  \item They correct a minor point that does not affect the outcome.
  \item Titles, precedents and who-decides matter more to them than what is decided.
  \item Praise from the right person changes their position within a conversation.
  \item They will not be seen to lose, even where losing costs nothing.
  \item Your success in a shared domain produces a change in temperature rather than in words.
\end{tells}
\begin{moves}
  \item Establish what rank the person believes they hold, and who they believe confers it.
  \item Address them at that rank --- slightly above it in public, exactly at it in private.
  \item Where you are junior, make your competence appear to reflect on them rather than to compete with them.
  \item Where you must oppose them, oppose the proposal and never the standing.
  \item Do not display an advantage you do not need to display.
\end{moves}
\begin{counters}
  \item Notice when a proposal is being resisted on grounds of who suggested it. That is status, not substance, and it can be answered by re-attributing the idea.
  \item Do not pay for praise with concessions. The two feel related and are not.
  \item Ask yourself which of your positions you would drop if nobody knew you held it. Those are the ones being held for rank.
  \item Where a superior is insecure, reduce visible competition rather than increasing visible competence --- the second reads as the first.
\end{counters}
\begin{cost}
Working the handle is easy and expensive. Status threats produce the longest-lived resentments in this book, and the reaction is disproportionate to the trigger in exactly the way described in T-03-01, so the operator frequently finds they have created an enemy over something they did not intend to do. Flattery also has a ceiling: once recognised it converts to contempt and destroys the relationship's remaining value.
\end{cost}

\sources{B3, B1}
\seesources
\seealso{T-03-01, T-04-01, T-09-01, T-18-03}
